% Options for packages loaded elsewhere
% Options for packages loaded elsewhere
\PassOptionsToPackage{unicode}{hyperref}
\PassOptionsToPackage{hyphens}{url}
%
\documentclass[
  11pt,
  ignorenonframetext,
  aspectratio=169,
]{beamer}
\newif\ifbibliography
\usepackage{pgfpages}
\setbeamertemplate{caption}[numbered]
\setbeamertemplate{caption label separator}{: }
\setbeamercolor{caption name}{fg=normal text.fg}
\beamertemplatenavigationsymbolsempty
% remove section numbering
\setbeamertemplate{part page}{
  \centering
  \begin{beamercolorbox}[sep=16pt,center]{part title}
    \usebeamerfont{part title}\insertpart\par
  \end{beamercolorbox}
}
\setbeamertemplate{section page}{
  \centering
  \begin{beamercolorbox}[sep=12pt,center]{section title}
    \usebeamerfont{section title}\insertsection\par
  \end{beamercolorbox}
}
\setbeamertemplate{subsection page}{
  \centering
  \begin{beamercolorbox}[sep=8pt,center]{subsection title}
    \usebeamerfont{subsection title}\insertsubsection\par
  \end{beamercolorbox}
}
% Prevent slide breaks in the middle of a paragraph
\widowpenalties 1 10000
\raggedbottom
\AtBeginPart{
  \frame{\partpage}
}
\AtBeginSection{
  \ifbibliography
  \else
    \frame{\sectionpage}
  \fi
}
\AtBeginSubsection{
  \frame{\subsectionpage}
}
\usepackage{iftex}
\ifPDFTeX
  \usepackage[T1]{fontenc}
  \usepackage[utf8]{inputenc}
  \usepackage{textcomp} % provide euro and other symbols
\else % if luatex or xetex
  \usepackage{unicode-math} % this also loads fontspec
  \defaultfontfeatures{Scale=MatchLowercase}
  \defaultfontfeatures[\rmfamily]{Ligatures=TeX,Scale=1}
\fi
\usepackage{lmodern}

\ifPDFTeX\else
  % xetex/luatex font selection
\fi
% Use upquote if available, for straight quotes in verbatim environments
\IfFileExists{upquote.sty}{\usepackage{upquote}}{}
\IfFileExists{microtype.sty}{% use microtype if available
  \usepackage[]{microtype}
  \UseMicrotypeSet[protrusion]{basicmath} % disable protrusion for tt fonts
}{}
\makeatletter
\@ifundefined{KOMAClassName}{% if non-KOMA class
  \IfFileExists{parskip.sty}{%
    \usepackage{parskip}
  }{% else
    \setlength{\parindent}{0pt}
    \setlength{\parskip}{6pt plus 2pt minus 1pt}}
}{% if KOMA class
  \KOMAoptions{parskip=half}}
\makeatother

\usepackage{color}
\usepackage{fancyvrb}
\newcommand{\VerbBar}{|}
\newcommand{\VERB}{\Verb[commandchars=\\\{\}]}
\DefineVerbatimEnvironment{Highlighting}{Verbatim}{commandchars=\\\{\}}
% Add ',fontsize=\small' for more characters per line
\usepackage{framed}
\definecolor{shadecolor}{RGB}{241,243,245}
\newenvironment{Shaded}{\begin{snugshade}}{\end{snugshade}}
\newcommand{\AlertTok}[1]{\textcolor[rgb]{0.68,0.00,0.00}{#1}}
\newcommand{\AnnotationTok}[1]{\textcolor[rgb]{0.37,0.37,0.37}{#1}}
\newcommand{\AttributeTok}[1]{\textcolor[rgb]{0.40,0.45,0.13}{#1}}
\newcommand{\BaseNTok}[1]{\textcolor[rgb]{0.68,0.00,0.00}{#1}}
\newcommand{\BuiltInTok}[1]{\textcolor[rgb]{0.00,0.23,0.31}{#1}}
\newcommand{\CharTok}[1]{\textcolor[rgb]{0.13,0.47,0.30}{#1}}
\newcommand{\CommentTok}[1]{\textcolor[rgb]{0.37,0.37,0.37}{#1}}
\newcommand{\CommentVarTok}[1]{\textcolor[rgb]{0.37,0.37,0.37}{\textit{#1}}}
\newcommand{\ConstantTok}[1]{\textcolor[rgb]{0.56,0.35,0.01}{#1}}
\newcommand{\ControlFlowTok}[1]{\textcolor[rgb]{0.00,0.23,0.31}{\textbf{#1}}}
\newcommand{\DataTypeTok}[1]{\textcolor[rgb]{0.68,0.00,0.00}{#1}}
\newcommand{\DecValTok}[1]{\textcolor[rgb]{0.68,0.00,0.00}{#1}}
\newcommand{\DocumentationTok}[1]{\textcolor[rgb]{0.37,0.37,0.37}{\textit{#1}}}
\newcommand{\ErrorTok}[1]{\textcolor[rgb]{0.68,0.00,0.00}{#1}}
\newcommand{\ExtensionTok}[1]{\textcolor[rgb]{0.00,0.23,0.31}{#1}}
\newcommand{\FloatTok}[1]{\textcolor[rgb]{0.68,0.00,0.00}{#1}}
\newcommand{\FunctionTok}[1]{\textcolor[rgb]{0.28,0.35,0.67}{#1}}
\newcommand{\ImportTok}[1]{\textcolor[rgb]{0.00,0.46,0.62}{#1}}
\newcommand{\InformationTok}[1]{\textcolor[rgb]{0.37,0.37,0.37}{#1}}
\newcommand{\KeywordTok}[1]{\textcolor[rgb]{0.00,0.23,0.31}{\textbf{#1}}}
\newcommand{\NormalTok}[1]{\textcolor[rgb]{0.00,0.23,0.31}{#1}}
\newcommand{\OperatorTok}[1]{\textcolor[rgb]{0.37,0.37,0.37}{#1}}
\newcommand{\OtherTok}[1]{\textcolor[rgb]{0.00,0.23,0.31}{#1}}
\newcommand{\PreprocessorTok}[1]{\textcolor[rgb]{0.68,0.00,0.00}{#1}}
\newcommand{\RegionMarkerTok}[1]{\textcolor[rgb]{0.00,0.23,0.31}{#1}}
\newcommand{\SpecialCharTok}[1]{\textcolor[rgb]{0.37,0.37,0.37}{#1}}
\newcommand{\SpecialStringTok}[1]{\textcolor[rgb]{0.13,0.47,0.30}{#1}}
\newcommand{\StringTok}[1]{\textcolor[rgb]{0.13,0.47,0.30}{#1}}
\newcommand{\VariableTok}[1]{\textcolor[rgb]{0.07,0.07,0.07}{#1}}
\newcommand{\VerbatimStringTok}[1]{\textcolor[rgb]{0.13,0.47,0.30}{#1}}
\newcommand{\WarningTok}[1]{\textcolor[rgb]{0.37,0.37,0.37}{\textit{#1}}}

\usepackage{longtable,booktabs,array}
\newcounter{none} % for unnumbered tables
\usepackage{calc} % for calculating minipage widths
\usepackage{caption}
% Make caption package work with longtable
\makeatletter
\def\fnum@table{\tablename~\thetable}
\makeatother
\usepackage{graphicx}
\makeatletter
\newsavebox\pandoc@box
\newcommand*\pandocbounded[1]{% scales image to fit in text height/width
  \sbox\pandoc@box{#1}%
  \Gscale@div\@tempa{\textheight}{\dimexpr\ht\pandoc@box+\dp\pandoc@box\relax}%
  \Gscale@div\@tempb{\linewidth}{\wd\pandoc@box}%
  \ifdim\@tempb\p@<\@tempa\p@\let\@tempa\@tempb\fi% select the smaller of both
  \ifdim\@tempa\p@<\p@\scalebox{\@tempa}{\usebox\pandoc@box}%
  \else\usebox{\pandoc@box}%
  \fi%
}
% Set default figure placement to htbp
\def\fps@figure{htbp}
\makeatother





\setlength{\emergencystretch}{3em} % prevent overfull lines

\providecommand{\tightlist}{%
  \setlength{\itemsep}{0pt}\setlength{\parskip}{0pt}}



 


\setbeamertemplate{navigation symbols}{}
\setbeamertemplate{footline}[frame number]
\setbeamertemplate{itemize items}[circle]
\setbeamertemplate{itemize subitem}[triangle]
\setbeamercolor{frametitle}{fg=black}
\setbeamerfont{frametitle}{series=\bfseries}
\newcommand{\indep}{\perp\!\!\!\perp}
\usepackage{booktabs}
\usepackage{tikz}
\usetikzlibrary{arrows.meta,positioning,calc}
\usepackage{fancyvrb}
\fvset{fontsize=\scriptsize}
\makeatletter
\def\verbatim@font{\scriptsize\ttfamily}
\makeatother
\makeatletter
\@ifpackageloaded{caption}{}{\usepackage{caption}}
\AtBeginDocument{%
\ifdefined\contentsname
  \renewcommand*\contentsname{Table of contents}
\else
  \newcommand\contentsname{Table of contents}
\fi
\ifdefined\listfigurename
  \renewcommand*\listfigurename{List of Figures}
\else
  \newcommand\listfigurename{List of Figures}
\fi
\ifdefined\listtablename
  \renewcommand*\listtablename{List of Tables}
\else
  \newcommand\listtablename{List of Tables}
\fi
\ifdefined\figurename
  \renewcommand*\figurename{Figure}
\else
  \newcommand\figurename{Figure}
\fi
\ifdefined\tablename
  \renewcommand*\tablename{Table}
\else
  \newcommand\tablename{Table}
\fi
}
\@ifpackageloaded{float}{}{\usepackage{float}}
\floatstyle{ruled}
\@ifundefined{c@chapter}{\newfloat{codelisting}{h}{lop}}{\newfloat{codelisting}{h}{lop}[chapter]}
\floatname{codelisting}{Listing}
\newcommand*\listoflistings{\listof{codelisting}{List of Listings}}
\makeatother
\makeatletter
\makeatother
\makeatletter
\@ifpackageloaded{caption}{}{\usepackage{caption}}
\@ifpackageloaded{subcaption}{}{\usepackage{subcaption}}
\makeatother

\usepackage{bookmark}
\IfFileExists{xurl.sty}{\usepackage{xurl}}{} % add URL line breaks if available
\urlstyle{same}
\hypersetup{
  pdftitle={Regression Discontinuity},
  pdfauthor={Professor Ji-Woong Chung},
  hidelinks,
  pdfcreator={LaTeX via pandoc}}


\title{\texorpdfstring{Regression
Discontinuity}{Regression Discontinuity}}
\subtitle{\texorpdfstring{BUSS975: Causal Inference in Financial
Research}{BUSS975: Causal Inference in Financial Research}}
\author{Professor Ji-Woong Chung}
\date{}
\institute{Korea University}

\begin{document}
\frame{\titlepage}


\section{6.1 Bird's-Eye View of RD}\label{birds-eye-view-of-rd}

\begin{frame}{Where we left off}
\protect\phantomsection\label{where-we-left-off}
\begin{itemize}
\tightlist
\item
  Lecture 3 ended on a trade. \textbf{Matching} needs common support;
  \textbf{regression} needs a correctly specified outcome model and will
  extrapolate where the data are silent.
\item
  In regression discontinuity, at any value of the running variable, you
  observe treated units \textbf{or} control units. Overlap is exactly
  zero.
\item
  So extrapolation is the entire design. The question becomes: \emph{how
  short a distance can we extrapolate, and how honestly?}
\item
  The answer --- extrapolate to a \textbf{point}, from both sides --- is
  what makes RD a credible observational design we have.
\end{itemize}
\end{frame}

\begin{frame}{The intuition: drinking at 21}
\protect\phantomsection\label{the-intuition-drinking-at-21}
\pandocbounded{\includegraphics[keepaspectratio]{04_regression_discontinuity_files/figure-beamer/unnamed-chunk-1-1.pdf}}

\vspace{-0.3em}
\footnotesize

Drinking rises at 21 because that is when it becomes legal.
Motor-vehicle mortality falls steadily with age, then jumps back up at
exactly the same point.
\end{frame}

\begin{frame}{What makes it an RD figure}
\protect\phantomsection\label{what-makes-it-an-rd-figure}
Three elements, and every RD figure has all three:

\begin{enumerate}
\tightlist
\item
  \textbf{The dots.} Average outcomes within narrow bins of the running
  variable --- here, one dot per month of age. Not raw observations.
\item
  \textbf{The lines.} Fitted values from a regression, estimated
  \textbf{separately on each side} of the cutoff. The curvature comes
  from a polynomial in age.
\item
  \textbf{The break.} That gap is the whole estimate.
\end{enumerate}

\vspace{0.3em}

\begin{itemize}
\tightlist
\item
  The break exists only because two different models were fit. Nothing
  forces them to meet --- so if they do meet, there is no effect.
\item
  \textbf{The magnitude.} Roughly \(5\) extra deaths per \(100{,}000\)
  at age 21. With about \(4.3\) million 21-year-olds in the US, that is
  \(5 \times 43 \approx \mathbf{215}\) additional deaths a year
  attributable to legal drinking.
\end{itemize}
\end{frame}

\begin{frame}{The jargon}
\protect\phantomsection\label{the-jargon}
\begin{itemize}
\tightlist
\item
  \textbf{Running variable} (\(X\)): the score that assigns treatment.
  Age, a test score, a credit score, a vote share, a BAC reading.
  Ideally continuous; a fine enough grid will do.
\item
  \textbf{Cutoff} or \textbf{threshold} (\(c_0\)): the known value where
  assignment changes.
\item
  \textbf{Trends}: how the \emph{outcome} moves with the running
  variable --- \textbf{not} time trends. In the mortality figure, deaths
  fall with age; that is an outcome trend.
\item
  \textbf{Discontinuity} or \textbf{jump}: the break in the conditional
  mean at \(c_0\). That break is the estimate.
\item
  The design is named for the discontinuity, but it lives or dies on the
  \textbf{trends}. Strong trends make effects hard to detect and easy to
  manufacture.
\end{itemize}
\end{frame}

\begin{frame}{Where do running variables come from?}
\protect\phantomsection\label{where-do-running-variables-come-from}
\begin{itemize}
\tightlist
\item
  They sound exotic. They are everywhere, because \textbf{humans write
  rules}, and rules need thresholds.
\item
  \textbf{Governments}: Medicare at 65, the drinking age at 21, adult
  criminal liability at 18, income thresholds for transfer programs.
\item
  \textbf{Firms and markets}: credit-score cutoffs in lending, index
  membership rules, covenant thresholds, EPS targets in compensation
  contracts, a \(50\%\) vote share.
\item
  \textbf{Algorithms}: surge pricing, credit limits, automated
  underwriting.
\item
  Finding one requires \textbf{institutional knowledge}, not
  econometrics. Knowing the rule exists, and exactly how it is applied,
  is the hard part.
\end{itemize}
\end{frame}

\begin{frame}{And RD needs a great deal of data}
\protect\phantomsection\label{and-rd-needs-a-great-deal-of-data}
\begin{itemize}
\tightlist
\item
  The estimate comes from observations \textbf{near the cutoff}.
  Everything else only helps pin down the two limits.
\item
  So the binding constraint is not the number of rows in your dataset.
  It is the number of rows in a neighbourhood of one point.
\item
  Strong trends make this worse: the more steeply the outcome moves with
  \(X\), the more data you need to separate the jump from the slope.
\item
  This is why RD studies are overwhelmingly built on
  \textbf{administrative records} --- birth registries, school records,
  court records, credit bureaus.
\item
  It is also why RD in finance is often underpowered. We return to this.
\end{itemize}
\end{frame}

\section{6.2 Identification
Assumptions}\label{identification-assumptions}

\begin{frame}{Sharp and fuzzy}
\protect\phantomsection\label{sharp-and-fuzzy}
\begin{center}
\small
\begin{tabular}{lll}
\toprule
 & Sharp & Fuzzy \\
\midrule
$\Pr(D = 1 \mid X)$ at $c_0$ & jumps $0 \to 1$ & jumps, but not to $1$ \\
Treatment is       & a deterministic function of $X$ & influenced by $X$ \\
The cutoff is      & the assignment rule & an \emph{instrument} \\
Example            & Medicare at $65$ & Israeli class-size rules \\
\bottomrule
\end{tabular}
\end{center}

\begin{itemize}
\tightlist
\item
  We do the \textbf{sharp} design here. The fuzzy design is an
  instrumental-variables problem, so it waits for the IV lecture.
\item
  Many finance applications are fuzzy. A shareholder proposal that
  passes with \(51\%\) is not always implemented.
\end{itemize}
\end{frame}

\begin{frame}{The sharp design}
\protect\phantomsection\label{the-sharp-design}
In a sharp RD, treatment status is a \textbf{deterministic,
discontinuous} function of the running variable: \[
D_i = \begin{cases} 1 & \text{if } X_i \geq c_0 \\ 0 & \text{if } X_i < c_0 \end{cases}
\]

\begin{itemize}
\tightlist
\item
  Know \(X_i\) and you know \(D_i\) with certainty. Know someone's age
  and their country, and you know whether they may legally drink.
\item
  Contrast the fuzzy design, where crossing \(c_0\) raises \(\Pr(D=1)\)
  without driving it to one.
\end{itemize}
\end{frame}

\begin{frame}{The key feature: no common support}
\protect\phantomsection\label{the-key-feature-no-common-support}
\begin{itemize}
\tightlist
\item
  This is a selection-on-observables design: condition on \(X\) and
  treatment is fully determined, so \textbf{unconfoundedness holds
  trivially}.
\item
  But the other assumption of Lecture 3 fails completely. At any value
  of \(X\) you see treated units or controls, never both: \[
  e(X) = \Pr(D = 1 \mid X) \in \{0, 1\} \quad \text{everywhere.}
  \]
\item
  No 19-year-old may legally drink; every 21-year-old may. There is no
  overlap to exploit.
\item
  \textbf{Every tool from last week is unusable.} You cannot match,
  stratify, or weight --- no unit has a counterpart at its own value of
  \(X\).
\item
  What replaces overlap is a \textbf{smoothness} assumption, plus a
  limiting argument at a single point.
\end{itemize}
\end{frame}

\begin{frame}{The target parameter}
\protect\phantomsection\label{the-target-parameter}
\pandocbounded{\includegraphics[keepaspectratio]{04_regression_discontinuity_files/figure-beamer/unnamed-chunk-2-1.pdf}}

\vspace{-0.3em}
\footnotesize

Left: parallel potential outcomes, so the vertical gap is the same at
every age and the number at the cutoff describes everyone. Right: the
functions curve and \textbf{cross at 23} --- the effect is \(5.5\) at
age 19, \(2.35\) at the cutoff, zero at 23, and \emph{negative} beyond.
RD returns the gap \textbf{at the dashed line} in both panels, and
nothing else.
\end{frame}

\begin{frame}{The estimand is a limit}
\protect\phantomsection\label{the-estimand-is-a-limit}
\[
\tau \;=\; \underbrace{\lim_{x \downarrow c_0} E[Y_i \mid X_i = x]}_{\text{from the right}} \;-\; \underbrace{\lim_{x \uparrow c_0} E[Y_i \mid X_i = x]}_{\text{from the left}}
\;=\; E[Y^1_i - Y^0_i \mid X_i = c_0]
\]

\begin{itemize}
\tightlist
\item
  It is only \emph{in the limit} that overlap is restored. Approaching
  \(c_0\) from either side, the two groups become arbitrarily similar in
  \(X\).
\item
  This is an \(ATE\), but for the subpopulation \(\{X = c_0\}\) ---
  which is why it is often called a \textbf{local} average treatment
  effect.
\item
  Note what it is \textbf{not}: not the \(ATE\), not the \(ATT\). It is
  an average at one point of one covariate.
\item
  Data from a wide range of \(X\) are used only to estimate those two
  limits precisely. The estimand still lives at \(c_0\).
\end{itemize}
\end{frame}

\begin{frame}{Smoothness: the assumption that does the work}
\protect\phantomsection\label{smoothness-the-assumption-that-does-the-work}
\begin{itemize}
\tightlist
\item
  \textbf{Hahn, Todd and Van der Klaauw (2001)} proved the effect at the
  cutoff is identified if the conditional expectations of \textbf{both}
  potential outcomes are continuous at \(c_0\): \[
  E[Y^0_i \mid X_i = x] \quad\text{and}\quad E[Y^1_i \mid X_i = x] \quad\text{are continuous at } x = c_0 .
  \]
\item
  This is the RD version of ``nature does not jump.'' The null is
  continuity; a jump is evidence that \emph{something} happened, and the
  design's job is to make treatment the only candidate.
\end{itemize}
\end{frame}

\begin{frame}{Smoothness, in a simulation}
\protect\phantomsection\label{smoothness-in-a-simulation}
\pandocbounded{\includegraphics[keepaspectratio]{04_regression_discontinuity_files/figure-beamer/unnamed-chunk-3-1.pdf}}

\vspace{-0.3em}
\footnotesize

\(Y^1 = 100 + 1.25X + \varepsilon\) and
\(Y^0 = 50 + 0.75X + \varepsilon\). Neither jumps at \(X = 250\) ---
\textbf{nothing happens at the cutoff} in the potential outcomes. Apply
the switching equation and the \emph{observed} outcome jumps anyway, by
exactly the vertical distance between the two functions there. That
distance is \(\tau\).
\end{frame}

\begin{frame}{The price of the design}
\protect\phantomsection\label{the-price-of-the-design}
\begin{itemize}
\tightlist
\item
  \textbf{You get} an effect that is credible, visible, and hard to fake
  --- identified by a rule you can read in a statute or a contract.
\item
  \textbf{You pay} by learning about one point. If the policy question
  is ``what would this do to everyone?'', RD does not answer it.
\item
  The gap between what RD identifies and what a policymaker wants is
  often large. A study of a drinking age of 21 says nothing about a
  drinking age of 18.
\item
  \textbf{Extrapolating away from the cutoff} requires extra assumptions
  (Angrist and Rokkanen 2015). Most papers do not, and should say so.
\item
  State the estimand as what it is: \emph{``the average effect for units
  at the threshold.''} Not ``the effect.''
\end{itemize}
\end{frame}

\section{6.3 Estimation}\label{estimation}

\begin{frame}{Two steps}
\protect\phantomsection\label{two-steps}
\textbf{Step 1. Recentre the running variable} at the cutoff: \[
\tilde X_i = X_i - c_0 .
\] Once centred, the running variable is zero at the cutoff, so the
coefficient on \(D\) is the jump \textbf{there} rather than at
\(X = 0\).

\textbf{Step 2. Regress the outcome on treatment, the running variable,
and their interaction:} \[
Y_i = \alpha + \tau D_i + \beta_1 \tilde X_i + \beta_2\, D_i \cdot \tilde X_i + \varepsilon_i .
\]

\begin{itemize}
\tightlist
\item
  \(\beta_1\) is the slope \textbf{below} the cutoff,
  \(\beta_1 + \beta_2\) the slope \textbf{above}, and \(\hat\tau\)
  estimates the jump.
\item
  The interaction is not optional. It lets the two expected
  potential-outcome functions have different slopes --- which is to say,
  it allows heterogeneous treatment effects. Drop it and any difference
  in trend is loaded onto \(\hat\tau\).
\end{itemize}
\end{frame}

\begin{frame}{Where the specification comes from}
\protect\phantomsection\label{where-the-specification-comes-from}
Model each potential outcome as a \(p\)th-order polynomial in the
centred running variable: \[
E[Y^0 \mid X] = \mu + \beta_{01}\tilde X + \cdots + \beta_{0p}\tilde X^{\,p}, \qquad
E[Y^1 \mid X] = \mu + \tau + \beta_{11}\tilde X + \cdots + \beta_{1p}\tilde X^{\,p}
\]

Apply the \textbf{switching equation},
\(Y_i = D_i Y^1_i + (1 - D_i)Y^0_i = Y^0_i + D_i(Y^1_i - Y^0_i)\), and
take expectations given \(X\): \[
E[Y \mid X] = E[Y^0 \mid X] + D\big(E[Y^1 \mid X] - E[Y^0 \mid X]\big)
\]

which is the regression you actually run: \[
Y_i = \mu + \beta_{01}\tilde X_i + \cdots + \beta_{0p}\tilde X_i^{\,p}
\;+\; \tau D_i + \gamma_1 D_i \tilde X_i + \cdots + \gamma_p D_i \tilde X_i^{\,p} + \varepsilon_i
\]

\begin{itemize}
\tightlist
\item
  with \(\gamma_k = \beta_{1k} - \beta_{0k}\). Interacting \(D\) with
  \textbf{every} polynomial term is what lets the two sides take
  genuinely different shapes.
\item
  The treatment effect \textbf{at the cutoff} is \(\tau\).
\end{itemize}
\end{frame}

\begin{frame}{Get the shape wrong and the answer is nonsense}
\protect\phantomsection\label{get-the-shape-wrong-and-the-answer-is-nonsense}
\pandocbounded{\includegraphics[keepaspectratio]{04_regression_discontinuity_files/figure-beamer/unnamed-chunk-4-1.pdf}}

\vspace{-0.3em}
\footnotesize

The truth is a \textbf{cubic} in the running variable and the treatment
effect is exactly \textbf{zero} --- the two potential outcomes are
identical, so they touch everywhere. Straight lines cannot bend where
the data bend, so they arrive at the cutoff at the wrong heights and
manufacture a jump out of nothing.
\end{frame}

\begin{frame}{The damage, quantified}
\protect\phantomsection\label{the-damage-quantified}
\begin{center}\small
\begin{tabular}{lrrr}
\toprule
Fitted polynomial & $X$ & $X, X^2$ & $X, X^2, X^3$ \\
\midrule
$\hat\tau$ & $-170,164$ & $59,136$ & $40$ \\
\bottomrule
\end{tabular}
\end{center}

\begin{itemize}
\tightlist
\item
  \textbf{The true effect is zero.} Underfit by two orders and you get
  about \(-170{,}000\); by one, about \(+59{,}000\). Get the order right
  and the estimate collapses into noise.
\item
  These are not merely biased --- they are wrong by \emph{five orders of
  magnitude}, and they would be reported as \textbf{highly statistically
  significant}.
\item
  The mechanism is the extrapolation. \(\hat\tau\) is the gap between
  two fitted values \textbf{at a boundary}, and boundary predictions are
  exactly where a misspecified polynomial errs most.
\end{itemize}
\end{frame}

\begin{frame}{But do not solve it with more polynomials}
\protect\phantomsection\label{but-do-not-solve-it-with-more-polynomials}
\begin{itemize}
\tightlist
\item
  The obvious fix --- keep raising \(p\) --- is the wrong one.
  \textbf{Gelman and Imbens (2019)} give three reasons not to go above
  quadratic:

  \begin{enumerate}
  \tightlist
  \item
    High-order polynomials put \textbf{large, erratic weights} on
    observations far from the cutoff, so distant data drive an estimate
    that is supposed to be local.
  \item
    Estimates are \textbf{highly sensitive} to the order chosen, with no
    principled way to pick it.
  \item
    The resulting confidence intervals have \textbf{poor coverage}.
  \end{enumerate}
\item
  Their recommendation, now standard: \textbf{local linear or local
  quadratic}, fitted only near the cutoff.
\item
  Which reframes the problem. Instead of ``what shape fits the whole
  range?'', ask ``how small a neighbourhood can I afford, and fit a line
  inside it.''
\end{itemize}
\end{frame}

\begin{frame}{The boundary problem with ``rectangular kernel''}
\protect\phantomsection\label{the-boundary-problem-with-rectangular-kernel}
\pandocbounded{\includegraphics[keepaspectratio]{04_regression_discontinuity_files/figure-beamer/unnamed-chunk-6-1.pdf}}

\vspace{-0.4em}
\footnotesize

The true effect is \(\overline{AB} = 2.4\). But a bin mean sits at the
bin's \textbf{centre}, not at its edge, so the left limit is read off at
\(A'\) and the right at \(B'\): the estimate is
\(\overline{A'B'} = 3.5\), too large by \(46\%\). \textbf{The trends
cause the bias} --- flatten the two lines and \(A' \to A\),
\(B' \to B\).
\end{frame}

\begin{frame}{Kernel and bandwidth}
\protect\phantomsection\label{kernel-and-bandwidth}
\begin{itemize}
\tightlist
\item
  The standard fix (Hahn, Todd and Van der Klaauw): fit a \textbf{line}
  inside the window rather than a mean, and let it extrapolate the short
  distance to the boundary. That is local linear regression, and it
  substantially reduces the bias.
\end{itemize}

\[
(\hat a, \hat b) \;=\; \arg\min_{a, b} \sum_{i=1}^n \big(Y_i - a - b(X_i - c_0)\big)^2\,
K\!\left(\frac{X_i - c_0}{h}\right) \mathbf{1}\{X_i > c_0\}
\]

\begin{itemize}
\tightlist
\item
  Think of kernel regression as a weighted regression restricted to a
  window. Two distinct choices, often confused:

  \begin{itemize}
  \tightlist
  \item
    The \textbf{kernel} \(K(\cdot)\) is the \emph{shape} of the window
    --- how weight declines with distance from the cutoff.
  \item
    The \textbf{bandwidth} \(h\) is the \emph{width} of the window ---
    how far from the cutoff you are willing to look.
  \end{itemize}
\item
  Averaging within a bin \textbf{is} a rectangular kernel: equal weight
  to every unit inside. That is exactly the biased case above.
\end{itemize}
\end{frame}

\begin{frame}{Three choices: \(h\), \(K(\cdot)\), and \(p\)}
\protect\phantomsection\label{three-choices-h-kcdot-and-p}
\begin{itemize}
\tightlist
\item
  \textbf{The bandwidth \(h\).} Smaller: fewer observations, less bias,
  more variance. Larger: more data, less variance, more bias if the
  included units depart from the true function. This choice matters
  most.
\item
  \textbf{The kernel \(K(\cdot)\).}

  \begin{itemize}
  \tightlist
  \item
    \emph{Rectangular}: equal weight inside the window, zero outside ---
    a plain window.
  \item
    \emph{Triangular}: weight declines linearly to zero at the edge.
  \item
    \emph{Epanechnikov}: parabolic weighting.
  \end{itemize}
\item
  \textbf{The polynomial order \(p\).} \(p = 0\) is a difference in
  means --- intuitive, but it \emph{is} the boundary problem. \(p > 1\)
  addresses it but risks overfitting and instability. \textbf{Local
  linear (\(p = 1\)) is generally preferred.}
\item
  Calonico, Cattaneo and Titiunik (2014) recommend the
  \textbf{triangular kernel with an MSE-optimal bandwidth}: maximum
  weight at the cutoff, zero at the window's edge.
\item
  There is no universally correct answer. Report what you chose, and
  show the answer does not turn on it.
\end{itemize}
\end{frame}

\begin{frame}{The three kernels, written out}
\protect\phantomsection\label{the-three-kernels-written-out}
Let \(u = (X_i - c_0)/h\), so that \(|u| \leq 1\) means ``inside the
window''.

\vspace{-0.3em}
\begin{center}
\small
\begin{tabular}{ll}
\toprule
Rectangular (uniform) & $K(u) = \tfrac{1}{2}\,\mathbf{1}\{|u| \leq 1\}$ \\
Triangular            & $K(u) = (1 - |u|)\,\mathbf{1}\{|u| \leq 1\}$ \\
Epanechnikov          & $K(u) = \tfrac{3}{4}(1 - u^2)\,\mathbf{1}\{|u| \leq 1\}$ \\
\bottomrule
\end{tabular}
\end{center}

\pandocbounded{\includegraphics[keepaspectratio]{04_regression_discontinuity_files/figure-beamer/unnamed-chunk-7-1.pdf}}

\vspace{-0.4em}
\footnotesize

\begin{itemize}
\tightlist
\item
  All three have \textbf{compact support} --- \(K(u) = 0\) once
  \(|u| > 1\).
\item
  The constants \(\tfrac12\) and \(\tfrac34\) make each integrate to
  one, and are \textbf{irrelevant here}: scaling every weight equally
  leaves the WLS solution unchanged.
\end{itemize}
\end{frame}

\begin{frame}[fragile]{The recipe}
\protect\phantomsection\label{the-recipe}
\begin{enumerate}
\tightlist
\item
  Choose a polynomial order \(p\) and a kernel \(K(\cdot)\).
\item
  Select a bandwidth \(h\).
\item
  \textbf{Above the cutoff}, run weighted least squares of \(Y\) on a
  constant and \(p\) terms in \((X_i - c_0)\), with weights
  \(K((X_i - c_0)/h)\). The intercept estimates
  \(\mu_+ = \lim_{x \downarrow c_0} E[Y \mid X = x]\): \[
  \hat Y_i = \hat\mu_+ + \hat\mu_{+,1}(X_i - c_0) + \cdots + \hat\mu_{+,p}(X_i - c_0)^p
  \]
\item
  \textbf{Repeat below the cutoff} for \(\hat\mu_-\).
\item
  The sharp RD point estimate is the difference of the two intercepts:
  \[
  \hat\tau = \hat\mu_+ - \hat\mu_- .
  \]
\end{enumerate}

\begin{itemize}
\tightlist
\item
  In practice this is one line: \texttt{rdrobust(y,\ x,\ c\ =\ 0)} in R,
  \texttt{rdrobust\ y\ x,\ c(0)} in Stata.
\end{itemize}
\end{frame}

\begin{frame}{What ``weighted least squares'' means in step 3}
\protect\phantomsection\label{what-weighted-least-squares-means-in-step-3}
Give unit \(i\) the weight \(w_i = K\!\big((X_i - c_0)/h\big)\) and
minimise the \textbf{weighted} sum of squares: \[
\hat\beta \;=\; \arg\min_{\beta} \sum_{i=1}^{n} w_i \big(Y_i - x_i'\beta\big)^2 ,
\qquad x_i = \big(1,\; \tilde X_i,\; \tilde X_i^2, \ldots, \tilde X_i^{\,p}\big)', \quad \tilde X_i = X_i - c_0
\] with the closed form \(\;\hat\beta = (X'WX)^{-1}X'WY\), where
\(W = \operatorname{diag}(w_1, \ldots, w_n)\).

\begin{itemize}
\tightlist
\item
  Units with \(|\tilde X_i| > h\) get \(w_i = 0\) and contribute nothing
  to either \(X'WX\) or \(X'WY\).
\item
  \textbf{Rectangular kernel \(\Rightarrow\) plain OLS on the
  subsample}, because every in-window weight is then equal.
\end{itemize}
\end{frame}

\begin{frame}{Choosing the bandwidth}
\protect\phantomsection\label{choosing-the-bandwidth}
Bias grows with \(h\), variance falls with it. Writing each as a
function of \(h\) and minimising: \[
\mathrm{MSE}(h) \;\approx\; \underbrace{B^2 h^{2(p+1)}}_{\text{bias}^2} + \underbrace{\frac{V}{nh}}_{\text{variance}}
\qquad\Longrightarrow\qquad
h_{MSE} \;\propto\; \underbrace{V^{\frac{1}{2p+3}}}_{\text{noise}}\;
\underbrace{B^{-\frac{2}{2p+3}}}_{\text{curvature}}\;
\underbrace{n^{-\frac{1}{2p+3}}}_{\text{sample size}}
\]

\begin{itemize}
\tightlist
\item
  \textbf{Noise \(V\): positive exponent.} Noisier data \(\Rightarrow\)
  a \textbf{wider} window, to average the noise away. If \(V \to 0\),
  then \(h \to 0\) --- use only points at the
  cutoff.\footnote{The Remix* states the $n$ and noise directions the other way round. }
\item
  \textbf{Curvature \(B\): negative.} More curvature \(\Rightarrow\) a
  \textbf{narrower} window. If \(B \to 0\) (a truly linear CEF),
  \(h \to \infty\) --- just run OLS on everything.
\item
  \textbf{Sample size \(n\): negative.} More data \(\Rightarrow\) a
  \textbf{narrower} window. The extra data buys \emph{less bias}, not a
  wider window --- though the count inside it,
  \(nh \propto n^{(2p+2)/(2p+3)}\), still grows.
\item
  Imbens and Kalyanaraman (2012) gave the first data-driven rule;
  Calonico, Cattaneo and Titiunik (2014) is the standard. Bandwidths may
  differ by side.
\end{itemize}

\footnotesize
\vspace{-0.2em}

\emph{The Remix} states the \(n\) and noise directions the other way
round.
\end{frame}

\begin{frame}{Why the obvious confidence interval is wrong}
\protect\phantomsection\label{why-the-obvious-confidence-interval-is-wrong}
\begin{itemize}
\tightlist
\item
  The local polynomial is a \textbf{nonparametric approximation} to an
  unknown function. It does not fit perfectly inside the window, so
  \(\hat\tau\) carries an approximation bias \(B\) as well as sampling
  noise. Honest inference needs \[
  CI = \big[(\hat\tau - B) \pm 1.96 \sqrt{V}\big].
  \]
\item
  The conventional interval,
  \(CI_{us} = [\hat\tau \pm 1.96\sqrt{\hat V}]\), simply \textbf{assumes
  \(B = 0\)} --- it centres the interval on a biased point estimate.
\item
  The consequence is \textbf{over-rejection} of the null of no effect.
  Cattaneo, Jansson and Ma discourage reporting these, and they are not
  credible under even moderate bias.
\item
  Worse, there is an internal contradiction: the MSE-optimal bandwidth
  is chosen \emph{to trade bias against variance}, so at that bandwidth
  the bias is deliberately \textbf{not} negligible.
\end{itemize}
\end{frame}

\begin{frame}[fragile]{Three intervals}
\protect\phantomsection\label{three-intervals}
\begin{center}
\small
\begin{tabular}{lccl}
\toprule
Interval & Centred at & Standard error & Verdict \\
\midrule
Conventional, $CI_{us}$        & $\hat\tau$          & $\sqrt{\hat V}$      & do not use \\
Bias-corrected, $CI_{bc}$      & $\hat\tau - \hat B$ & $\sqrt{\hat V}$      & better, but ignores noise in $\hat B$ \\
Robust bias-corrected, $CI_{rbc}$ & $\hat\tau - \hat B$ & $\sqrt{\hat V_{bc}}$ & \textbf{report this} \\
\bottomrule
\end{tabular}
\end{center}

\begin{itemize}
\tightlist
\item
  Bias correction subtracts an estimate of \(B\). But \(\hat B\) is
  itself estimated, and pretending otherwise distorts coverage in finite
  samples.
\item
  \textbf{Robust} bias correction folds the variability of \(\hat B\)
  into the variance term. Intervals get somewhat wider for the same
  bandwidth --- that is the honest price of the correction.
\item
  This is what \texttt{rdrobust} reports by default. Report the robust
  row, and say that you did.
\end{itemize}
\end{frame}

\begin{frame}[fragile]{Two more things about standard errors}
\protect\phantomsection\label{two-more-things-about-standard-errors}
\begin{itemize}
\tightlist
\item
  \textbf{Do not cluster on the running variable.} Lee (2008) and Lee
  and Card (2008) recommended it and the practice spread through most
  subsequent RD studies. \textbf{Kolesár and Rothe (2018)} show it
  performs badly, with inflated Type I error. Use
  heteroskedasticity-robust errors, or ``honest'' intervals
  (\texttt{RDHonest}).
\item
  \textbf{Sparse data near the cutoff} is the recurring practical
  problem. When observations are thin, or the running variable is coarse
  with large gaps between values, robust intervals may fail to centre
  properly --- more often than practitioners realise.
\item
  \textbf{Randomization inference} (Cattaneo, Frandsen and Titiunik)
  offers an alternative: assume that within a \textbf{small window}
  around the cutoff, treatment is effectively randomly assigned, then
  analyse that window as an experiment.
\item
  It gives exact or approximate \(p\)-values under weaker assumptions
  than local polynomials, and does better when data are sparse. It also
  matches the experimental analogy RD keeps invoking.
\end{itemize}
\end{frame}

\section{6.4 Challenges to
Identification}\label{challenges-to-identification}

\begin{frame}{Five ways continuity fails}
\protect\phantomsection\label{five-ways-continuity-fails}
Continuity is the whole design. It breaks when:

\begin{enumerate}
\tightlist
\item
  The \textbf{assignment rule is known} in advance;
\item
  Agents are \textbf{interested} in adjusting their score;
\item
  Agents \textbf{have time} to adjust;
\item
  The \textbf{cutoff is endogenous} --- something else changes at
  \(c_0\) too;
\item
  There is \textbf{nonrandom heaping} along the running variable.
\end{enumerate}

\vspace{0.3em}

\begin{itemize}
\tightlist
\item
  Conditions 1--3 together are the recipe for \textbf{sorting}: retaking
  an exam, reporting income just below a threshold, delaying a birth.
  All three are needed --- an agent with an incentive \emph{and} an
  opportunity.
\item
  Condition 4 is different, and often overlooked. Turning 18 brings
  adult criminal penalties --- and also voting rights, high-school
  graduation, and much else. The cutoff is real but not clean, and the
  confound is not a variable you forgot to control for.
\item
  Condition 5 is mechanical rather than strategic, and gets its own
  treatment below.
\end{itemize}
\end{frame}

\begin{frame}{Sorting, and the two rooms}
\protect\phantomsection\label{sorting-and-the-two-rooms}
\begin{itemize}
\tightlist
\item
  Two rooms of patients waiting for a life-saving treatment. Room A will
  receive it; room B knowingly receives nothing. What do you do if you
  are in room B?
\item
  You stand up, open the door, and walk across the hall. The only thing
  that would keep you in room B is if crossing were impossible.
\item
  And if everyone crosses, room A ends up crowded and room B nearly
  empty. That imbalance is the \textbf{observable footprint} of sorting:
  you cannot see motives, but you can count how many ended up on each
  side.
\item
  Formally: with a desirable treatment and rule \(X \geq c_0\), agents
  choose \(X\) just above \(c_0\) whenever they are able. If their
  \emph{ability} to do so is related to their potential outcomes, the
  comparison at the cutoff is contaminated.
\item
  Behavioural theory makes a \textbf{testable prediction} --- a
  discontinuity in the \emph{density} of the running variable.
  Statistics built on optimising behaviour can take you further than
  statistics alone.
\end{itemize}
\end{frame}

\begin{frame}[fragile]{The McCrary density test}
\protect\phantomsection\label{the-mccrary-density-test}
\pandocbounded{\includegraphics[keepaspectratio]{04_regression_discontinuity_files/figure-beamer/unnamed-chunk-8-1.pdf}}

\vspace{-0.4em}

\begin{itemize}
\tightlist
\item
  Under the null of no manipulation the density of \(X\) is
  \textbf{continuous} at \(c_0\); under the alternative it bunches on
  the treated side.
\item
  Mechanically: partition the running variable into bins, count
  observations per bin, treat the counts as an outcome, fit a local
  linear regression on each side, and test whether the limits agree.
\item
  Use \texttt{rddensity} (Cattaneo, Jansson and Ma) --- local
  polynomials with less bias at the border than the original. \textbf{It
  needs a lot of data at \(c_0\)} to separate a real jump from noise.
\end{itemize}
\end{frame}

\begin{frame}{Covariate balance and placebo cutoffs}
\protect\phantomsection\label{covariate-balance-and-placebo-cutoffs}
\begin{itemize}
\tightlist
\item
  \textbf{Covariate balance.} Pre-treatment characteristics cannot be
  caused by treatment, so they must not jump at \(c_0\). Run the same RD
  with each covariate as the outcome and find nothing. Lee, Moretti and
  Butler (2004) is the template: real income, share with a high-school
  degree, share Black, share eligible to vote --- all smooth through the
  \(50\%\) Democratic vote share.
\item
  A jump in a covariate is serious. It means something \emph{other than}
  treatment changes at the cutoff --- and if an observed characteristic
  jumps, unobserved ones probably do too.
\item
  \textbf{Placebo cutoffs.} Imbens and Lemieux (2008): take one side of
  the discontinuity, use the median of the running variable there as a
  fake cutoff \(c_0'\), and test for a jump in the outcome. You do not
  want to find anything.
\item
  The identifying assumptions cannot be tested directly. What these do
  is supply \textbf{indirect evidence} that may be persuasive ---
  placebo logic, looking where there should be no effect.
\end{itemize}
\end{frame}

\begin{frame}{Nonrandom heaping}
\protect\phantomsection\label{nonrandom-heaping}
\begin{itemize}
\tightlist
\item
  \textbf{Almond, Doyle, Kowalski and Williams (2010)} wanted the
  returns to medical care --- hard, because doctors assign care by
  expected outcome, which is selection bias by construction. The Perfect
  Doctor problem again.
\item
  Their insight: babies below \(1{,}500\)g are classified ``very low
  birth weight'' and receive heightened attention. Using hospital
  records linked to mortality, they found one-year mortality about
  \textbf{one percentage point lower} just below the threshold, against
  a mean of \(5.5\%\). Large, and apparently far exceeding the cost.
\item
  They ran the McCrary test at \(1{,}500\)g and \textbf{found nothing}.
  Satisfied, they proceeded.
\item
  \textbf{Barreca, Guldi, Lindo and Waddell} then looked at the
  \emph{shape} of the whole distribution rather than at one
  discontinuity, and found \textbf{heaping}: excess mass at every round
  \(100\)g mark, from hospitals rounding to the nearest integer --- or,
  more troublingly, from staff rounding \emph{to} \(1{,}500\)g to
  trigger specialised care.
\item
  Nature does not produce piles of babies every hundred grams. Heaping
  is either rounding or manipulation.
\end{itemize}
\end{frame}

\begin{frame}{What passing the density test does \emph{not} prove}
\protect\phantomsection\label{what-passing-the-density-test-does-not-prove}
\begin{itemize}
\tightlist
\item
  The babies heaped exactly at \(1{,}500\)g had \textbf{unusually high
  mortality} --- outliers relative to units on \emph{both} sides.
\item
  The consequence, in Barreca et al.'s words: the heaping alone makes it
  look ``good'' to be \textbf{strictly below any} \(100\)g cutoff
  between \(1{,}000\) and \(3{,}000\)g. That is not a treatment effect;
  it is a composition effect.
\item
  The McCrary test asks one question --- \emph{is the density
  continuous?} --- and the answer was yes. The problem lived in a
  different feature of the same distribution.
\item
  \textbf{The lesson is the one from Lecture 3.} A test on one feature
  cannot clear a problem living in another. There, means passed while
  variances failed. Here, the density was continuous while the
  composition was not. \hfill \emph{(added)}
\item
  Plot the raw distribution. Do not rely on a test statistic alone.
\end{itemize}
\end{frame}

\begin{frame}{The donut-hole RD}
\protect\phantomsection\label{the-donut-hole-rd}
\begin{itemize}
\tightlist
\item
  If units \emph{at} the cutoff are suspect --- heaped, rounded, or
  manipulated --- but units nearby are fine, drop the suspect ones and
  re-estimate. That is the \textbf{donut hole}.
\item
  The justification is clean: RD compares means \emph{as we approach}
  the threshold, so the estimate should not be sensitive to observations
  exactly at it.
\item
  In Almond et al.'s data, dropping about \textbf{\(2\%\)} of
  observations cut the estimated one-year mortality effect roughly
  \textbf{in half}.
\item
  The costs are real:

  \begin{itemize}
  \tightlist
  \item
    You lose observations you often cannot afford to lose, so it may be
    infeasible when the heap is large relative to the sample.
  \item
    The estimand changes. It is now local to a population defined by a
    hole in its middle --- harder to describe, and further still from
    anything a policymaker asked about.
  \end{itemize}
\item
  Heaping is not fatal. But report both, with and without the donut, and
  let the reader see the difference.
\end{itemize}
\end{frame}

\section{6.5 Medicare and Universal
Healthcare}\label{medicare-and-universal-healthcare}

\begin{frame}{The question}
\protect\phantomsection\label{the-question}
\begin{itemize}
\tightlist
\item
  \textbf{Card, Dobkin and Maestas (2008)}: what does universal health
  insurance do to healthcare use? Medicare arrives sharply at age
  \textbf{65}.
\item
  The stakes are large. In 2014 Medicare was \(14\%\) of the US federal
  budget --- \(\$505\) billion.
\item
  Around \(20\%\) of non-elderly US adults were uninsured in 2005,
  mostly from lower-income families, nearly half Black or Hispanic.
  Unequal coverage is a leading explanation for disparities in
  healthcare use and health outcomes.
\item
  \textbf{Why observational comparisons fail.} Both the supply of and
  the demand for insurance depend on health status, so
  \(\mathrm{cov}(u, C) \neq 0\). Comparing the insured with the
  uninsured measures selection as much as treatment --- and even among
  the insured, copays and deductibles vary with health.
\item
  For the elderly the picture inverts: \textbf{under \(1\%\) are
  uninsured}, most on fee-for-service Medicare, and the transition
  happens at one known age.
\end{itemize}
\end{frame}

\begin{frame}{The empirical design}
\protect\phantomsection\label{the-empirical-design}
Healthcare use \(y\) depends on characteristics \(C^k\) of coverage
(copayments, generosity), for individual \(i\), group \(j\), age \(a\).
\textbf{\(C\) is endogenous}, so OLS does not identify \(\pi_k\): \[
y_{ija} = X_{ija}\beta + f_j(a) + \textstyle\sum_k \pi_k C^k_{ija} + u_{ija}
\]

Model two coverage dummies --- any coverage, and generous coverage ---
as jumping at 65, with \(D_a = \mathbf{1}\{a \geq 65\}\): \[
C^k_{ija} = X_{ija}\beta^k_j + g^k_j(a) + D_a\,\delta^k_j + v^k_{ija}, \qquad k = 1, 2
\]

Substituting, the jump in \(y\) at 65 is a \textbf{combination} of the
two coverage jumps:
\(\;\delta^y_j = \pi_1 \delta^1_j + \pi_2 \delta^2_j\).

\begin{itemize}
\tightlist
\item
  If \(f_j\), \(g^1_j\) and \(g^2_j\) are all \textbf{continuous at 65}
  --- the identifying assumption --- any discontinuity in \(y\) is due
  to insurance.
\item
  One group gives one equation in two unknowns. With \textbf{several
  groups}, regressing \(\delta^y_j\) on \(\delta^1_j\) and
  \(\delta^2_j\) across \(j\) separates \(\pi_1\) from \(\pi_2\).
\end{itemize}
\end{frame}

\begin{frame}{The data}
\protect\phantomsection\label{the-data}
\begin{itemize}
\tightlist
\item
  \textbf{NHIS, 1992--2003.} Reports birth year, birth month, and
  calendar quarter of interview --- enough to construct \textbf{age in
  quarters}. Someone reaching 65 in the interview quarter is coded 65
  and 0 quarters; assuming uniform interview dates, half of them are up
  to six weeks younger than 65 and half up to six weeks older.
\item
  Restricted to ages 55--75. Final sample: \textbf{\(160{,}821\)}
  observations.
\item
  \textbf{Hospital discharge records} for California, Florida and New
  York --- a complete census of discharges, with \textbf{age in months}
  at admission. Transfers dropped, ages 60--70 kept. Sample sizes
  \(4{,}017{,}325\), \(2{,}793{,}547\) and \(3{,}121{,}721\).
\item
  \textbf{The institution is genuinely sharp.} Coverage begins on the
  first day of the month you turn 65. Individuals are notified shortly
  before their birthday, and must enrol and choose whether to take Part
  B. Part A is free; Part B costs a modest premium.
\end{itemize}
\end{frame}

\begin{frame}{What jumps at 65: insurance}
\protect\phantomsection\label{what-jumps-at-65-insurance}
\begin{center}
\scriptsize
\begin{tabular}{lrrrrr}
\toprule
 & On Medicare & Any insurance & Private & 2+ forms & Managed care \\
\midrule
Overall sample        & $59.7$ & $9.5$  & $-2.9$ & $44.1$ & $-28.4$ \\
                      & $(4.1)$ & $(0.6)$ & $(1.1)$ & $(2.8)$ & $(2.1)$ \\[0.3em]
White, some college   & $68.4$ & $4.4$  & $-2.3$ & $55.1$ & $-40.1$ \\
                      & $(4.7)$ & $(0.5)$ & $(1.8)$ & $(4.0)$ & $(2.6)$ \\[0.3em]
Minority, HS dropout  & $44.5$ & $21.5$ & $-1.2$ & $19.4$ & $-8.3$ \\
                      & $(3.1)$ & $(2.1)$ & $(2.5)$ & $(1.9)$ & $(3.1)$ \\
\bottomrule
\end{tabular}
\end{center}

\footnotesize

Percentage-point discontinuities at 65, from quadratics in age fully
interacted with the post-65 dummy, plus controls for gender, education,
race, region and year. Pooled 1999--2003 NHIS. \normalsize

\begin{itemize}
\tightlist
\item
  The \textbf{first stage is enormous}: Medicare coverage jumps about
  \(60\) points, and holding two or more forms of coverage jumps \(44\).
\item
  Who \emph{gains} insurance and who merely \textbf{swaps} it differs
  sharply. Minority dropouts gain coverage (\(+21.5\)); white
  college-goers were already insured and mostly change its form
  (\(+4.4\), with managed care down \(40\) points).
\end{itemize}
\end{frame}

\begin{frame}{The confounder: what else happens at 65?}
\protect\phantomsection\label{the-confounder-what-else-happens-at-65}
\begin{itemize}
\tightlist
\item
  Continuity requires that \textbf{everything else} trends smoothly
  through 65 --- observed and unobserved. But 65 is also the traditional
  retirement age, and non-workers have more time to visit a doctor.
\item
  If employment fell discontinuously at 65, part of any jump in
  healthcare use would be leisure, not insurance.
\item
  So the authors bring in a \textbf{third dataset}, the March CPS
  1996--2004, and run the same RD with employment as the outcome. They
  find \textbf{no discontinuity at 65}.
\item
  This is the covariate-balance logic of §6.4, aimed at the one
  confounder that institutional knowledge flags as dangerous.
\item
  \textbf{The general lesson.} ``Does anything else change at the
  cutoff?'' is not a statistical question first. You answer it by
  knowing the institution, and then testing what that knowledge
  suggests. \hfill \emph{(added)}
\end{itemize}
\end{frame}

\begin{frame}{What jumps at 65: access and use}
\protect\phantomsection\label{what-jumps-at-65-access-and-use}
\begin{center}
\scriptsize
\begin{tabular}{lrrrr}
\toprule
 & Delayed care & Did not get care & Saw a doctor & Hospital stay \\
\midrule
Overall sample        & $-1.8$ & $-1.3$ & $1.3$ & $1.2$ \\
                      & $(0.4)$ & $(0.3)$ & $(0.7)$ & $(0.4)$ \\[0.3em]
White non-Hispanic    & $-1.6$ & $-1.2$ & $0.6$ & $1.3$ \\
                      & $(0.4)$ & $(0.3)$ & $(0.8)$ & $(0.5)$ \\[0.3em]
Hispanic              & $-4.9$ & $-3.8$ & $8.2$ & $11.8$ \\
                      & $(0.8)$ & $(0.7)$ & $(0.8)$ & $(1.6)$ \\
\bottomrule
\end{tabular}
\end{center}

\footnotesize

Percentage-point discontinuities at 65. First two columns from the
1997--2003 NHIS, last two from 1992--2003. \normalsize

\begin{itemize}
\tightlist
\item
  Overall the effects are \textbf{small but precisely estimated}: care
  delayed over cost falls \(1.8\) points, unmet need \(1.3\), doctor
  visits rise \(1.3\), hospital stays \(1.2\).
\item
  The averages hide the story. For \textbf{Hispanics} the effects are
  three to ten times larger --- an \(11.8\)-point rise in hospital
  stays. Coverage matters most for those who had least.
\end{itemize}
\end{frame}

\begin{frame}{And mortality}
\protect\phantomsection\label{and-mortality}
\begin{center}
\scriptsize
\begin{tabular}{lrrrrrr}
\toprule
Death rate within & 7 days & 14 days & 28 days & 90 days & 180 days & 365 days \\
\midrule
Quadratic + controls & $-1.0$ & $-0.8$ & $-0.9$ & $-0.9$ & $-0.8$ & $-0.7$ \\
                     & $(0.2)$ & $(0.2)$ & $(0.3)$ & $(0.3)$ & $(0.3)$ & $(0.4)$ \\[0.3em]
Cubic + controls     & $-0.7$ & $-0.7$ & $-0.6$ & $-0.9$ & $-0.9$ & $-0.4$ \\
                     & $(0.3)$ & $(0.2)$ & $(0.4)$ & $(0.4)$ & $(0.5)$ & $(0.5)$ \\[0.3em]
Local OLS, ad hoc bandwidths & $-0.8$ & $-0.8$ & $-0.8$ & $-0.9$ & $-1.1$ & $-0.8$ \\
                     & $(0.2)$ & $(0.2)$ & $(0.2)$ & $(0.2)$ & $(0.3)$ & $(0.3)$ \\
\bottomrule
\end{tabular}
\end{center}

\footnotesize

From Card, Dobkin and Maestas (2009), a follow-up study.
Percentage-point discontinuities at 65. The chapter's fourth row --- a
quadratic with no controls --- is omitted here. \normalsize

\begin{itemize}
\tightlist
\item
  Roughly a \textbf{one-percentage-point} reduction in death rates,
  stable from 7 days out to a year.
\item
  Note what the table is really doing: reporting the same estimate under
  a quadratic, a cubic, and local OLS. \textbf{Answers that barely move
  across specifications} are the point of the exercise --- compare the
  §6.3 simulation, where they moved by five orders of magnitude.
\end{itemize}
\end{frame}

\begin{frame}{What to take from Card, Dobkin and Maestas}
\protect\phantomsection\label{what-to-take-from-card-dobkin-and-maestas}
\begin{itemize}
\tightlist
\item
  \textbf{The design is textbook sharp.} A statutory age, a first stage
  of \(60\) percentage points, and an institution that gives individuals
  no way to move their birthday.
\item
  \textbf{The chain is explicit.} Coverage jumps \(\rightarrow\) access
  improves \(\rightarrow\) utilisation rises \(\rightarrow\) mortality
  falls slightly. Every link is estimated at the same cutoff, in the
  same design.
\item
  \textbf{The confounder was named before it was tested.} Retirement at
  65 is the obvious threat, and they went to a third dataset to rule it
  out.
\item
  \textbf{Heterogeneity is where the economics is.} On average Medicare
  mostly \emph{swaps} one form of coverage for another, and the average
  effect on access is small. Both averages conceal large gains for the
  previously uninsured.
\item
  Contrast with Hansen, coming next: there the estimate is large and the
  \emph{mechanism} stays ambiguous. Here the mechanism is clear and the
  average effect is small. \hfill \emph{(added)}
\end{itemize}
\end{frame}

\section{6.6 DUI and Recidivism: A Data
Exercise}\label{dui-and-recidivism-a-data-exercise}

\begin{frame}{Why deterrence is hard to identify}
\protect\phantomsection\label{why-deterrence-is-hard-to-identify}
\begin{itemize}
\tightlist
\item
  Becker (1968): crime is deterred either by raising the
  \textbf{probability of arrest} or by raising the \textbf{punishment}
  conditional on conviction. Both are policy levers, and the ``scared
  straight'' reading says the mere communication of a severe penalty may
  be enough.
\item
  The empirical problem is that harsher penalties do two things at once.
  They may \textbf{deter} --- offenders rationally reoffend less --- and
  they certainly \textbf{incapacitate}, by removing offenders from the
  population at risk.
\item
  Longer sentences are the leading example, and the main driver of US
  mass incarceration. Crime falls, but is that rational deterrence or
  simply that more people are behind bars? Aggregate data cannot
  separate the two.
\item
  Two designs that can: \textbf{Drago, Galbiati and Vertova (2009)}
  exploit an Italian collective pardon that left otherwise-identical
  offenders with different residual sentences. \textbf{Lee and McCrary
  (2017)} use exact birth dates and the jump in penalties at 18. Both
  find small deterrence effects.
\item
  \textbf{Hansen (2015)} found a third: something close to a coin flip
  inside a breathalyzer.
\end{itemize}
\end{frame}

\begin{frame}{The setting}
\protect\phantomsection\label{the-setting}
\begin{itemize}
\tightlist
\item
  In Washington State, a driver stopped on suspicion blows into a
  breathalyzer, which records \textbf{blood alcohol content}. At
  \(\mathbf{0.08}\) or above they are charged with DUI. Punishments
  \textbf{escalate with prior offences}, and a second threshold at
  \(0.15\) marks ``aggravated'' DUI --- we use \(0.08\) only.
\end{itemize}

\begin{center}
\scriptsize
\begin{tabular}{lrrrr}
\toprule
 & \multicolumn{2}{c}{1st offence} & \multicolumn{2}{c}{2nd offence} \\
\cmidrule(lr){2-3}\cmidrule(lr){4-5}
 & DUI & Agg. DUI & DUI & Agg. DUI \\
\midrule
Minimum fine        & \$865.50 & \$1{,}120.50 & \$1{,}120.50 & \$1{,}545.50 \\
Minimum jail time   & 24 hours & 48 hours & 30 days & 45 days \\
Licence suspension  & 90 days & 365 days & 2 years & 900 days \\
\bottomrule
\end{tabular}
\end{center}

\begin{itemize}
\tightlist
\item
  \textbf{The design:} compare drivers who blew \(0.079\) with drivers
  who blew \(0.081\). If those just above reoffend less, the escalated
  punishment deterred them.
\item
  \textbf{The outcome:} recidivism --- the same person stopped again and
  blowing above \(0.08\), measurable because the records carry
  persistent personal identifiers.
\item
  Our extract has \(214{,}558\) stops against Hansen's \(512{,}694\), so
  we \textbf{copy} his steps rather than replicate his numbers.
\end{itemize}
\end{frame}

\begin{frame}[fragile]{Continuous, or merely fine-grained?}
\protect\phantomsection\label{continuous-or-merely-fine-grained}
\begin{itemize}
\tightlist
\item
  Washington's breathalyzers reported to \(0.001\). That is precise ---
  but it is not continuous, and the distinction matters.
\item
  \textbf{Truly continuous} running variable: the probability that any
  two units share a value is zero.
\item
  \textbf{Multi-valued} running variable: the scale is coarsened, so
  many units share values. Nobody can sit \emph{between} two \(0.001\)
  bins.
\item
  Consequences, both practical:

  \begin{itemize}
  \tightlist
  \item
    Density tests ``can be a little funky'' when the running variable is
    not continuous --- there is mass \emph{at} points, not merely near
    them.
  \item
    \texttt{rdrobust} warns \emph{``Mass points detected in the running
    variable''} on these data, and \textbf{Kolesár and Rothe (2018)} was
    written for exactly this case.
  \end{itemize}
\item
  It rarely invalidates the design. But it is the kind of institutional
  detail that belongs in the paper rather than in a footnote.
  \hfill \emph{(added)}
\end{itemize}
\end{frame}

\begin{frame}[fragile]{Step 1: is the running variable manipulated?}
\protect\phantomsection\label{step-1-is-the-running-variable-manipulated}
\pandocbounded{\includegraphics[keepaspectratio]{04_regression_discontinuity_files/figure-beamer/unnamed-chunk-9-1.pdf}}

\vspace{-0.3em}
\footnotesize

No bunching at \(0.08\), and --- unlike the birth-weight data --- no
heaping at round numbers. \texttt{rddensity} gives \(T = -0.14\),
\(p = 0.89\). That is what the institution predicts. \textbf{Recall the
five conditions:} the rule is known and drivers would dearly like to
adjust, but the officer operates the machine and they have no
opportunity. Two of three conditions are not enough.
\end{frame}

\begin{frame}{Step 2: do covariates jump at the cutoff?}
\protect\phantomsection\label{step-2-do-covariates-jump-at-the-cutoff}
\begin{center}
\small
\begin{tabular}{lrrrr}
\toprule
 & white & male & age & accident \\
\midrule
Effect of DUI & $0.0057$ & $0.0062$ & $-0.1405$ & $-0.0034$ \\
              & $(0.0050)$ & $(0.0057)$ & $(0.1644)$ & $(0.0041)$ \\
Mean of outcome & $0.853$ & $0.792$ & $34.17$ & $0.101$ \\
$N$ & $89{,}967$ & $89{,}967$ & $89{,}967$ & $89{,}967$ \\
\bottomrule
\end{tabular}
\end{center}

\begin{itemize}
\tightlist
\item
  Hansen's specification:
  \(Y_i = \alpha + \tau\, \mathrm{DUI}_i + \beta_1 \mathrm{BAC}_i + \beta_2 \mathrm{BAC}_i \cdot \mathrm{DUI}_i + X_i\gamma + \varepsilon_i\),
  on \(0.03 \leq \mathrm{BAC} \leq 0.13\) --- a \textbf{rectangular}
  kernel with bandwidth \(0.05\) on each side, robust standard errors.
  Unweighted OLS \emph{is} a rectangular kernel.
\item
  Nothing is significant, and the point estimates are tiny against the
  means: \(0.6\) percentage points on an \(85\%\) base, a tenth of a
  year on a mean age of \(34\).
\item
  These are \textbf{precise nulls}, which is the useful kind. A null
  with enormous standard errors would tell you nothing at all.
\end{itemize}
\end{frame}

\begin{frame}{Step 2, in pictures}
\protect\phantomsection\label{step-2-in-pictures}
\pandocbounded{\includegraphics[keepaspectratio]{04_regression_discontinuity_files/figure-beamer/unnamed-chunk-10-1.pdf}}

\vspace{-0.3em}
\footnotesize

Means in \(0.002\)-wide BAC bins, with a line fitted on each side. The
lines \textbf{meet} at \(0.08\) in all four panels --- which is what the
regression table said, now visible. Note how much more convincing the
picture is than the coefficients were.
\end{frame}

\begin{frame}{Step 3: look at the outcome}
\protect\phantomsection\label{step-3-look-at-the-outcome}
\pandocbounded{\includegraphics[keepaspectratio]{04_regression_discontinuity_files/figure-beamer/unnamed-chunk-11-1.pdf}}

\vspace{-0.3em}
\footnotesize

Now the lines \textbf{do not} meet. Recidivism runs between \(8\%\) and
\(14\%\) --- not common, so the estimated effect will be large in
relative terms --- and it falls discontinuously for drivers who blew
just over the limit. Eyeballing the left panel: about \(11.5\%\) drops
to about \(9.75\%\), so roughly \(1.75\) points. The quadratic gives a
slightly smaller gap, nearer \(1.25\)--\(1.5\).
\end{frame}

\begin{frame}{Step 4: the regressions}
\protect\phantomsection\label{step-4-the-regressions}
\begin{center}
\small
\begin{tabular}{lcc}
\toprule
Outcome: recidivism & (1) & (2) \\
\midrule
Effect of DUI & $-0.0236$ & $-0.0143$ \\
              & $(0.0044)$ & $(0.0062)$ \\
\midrule
Polynomial & BAC & BAC, BAC$^2$ \\
Mean of outcome & $0.107$ & $0.107$ \\
$N$ & $89{,}967$ & $89{,}967$ \\
\bottomrule
\end{tabular}
\end{center}

\begin{itemize}
\tightlist
\item
  Same window and kernel as the balance table; column (2) interacts DUI
  with BAC \textbf{and} BAC\(^2\), modelling the potential-outcome
  functions more flexibly. Coefficients and standard errors reproduce
  Hansen's Table 6.7 exactly.
\item
  Column (1): a \(2.4\) percentage-point fall against a base of
  \(10.7\%\) --- a drop of more than a fifth, taking recidivism at the
  cutoff to about \(9.3\%\). Column (2) allows curvature and gives
  \(1.4\) points.
\item
  Both match the eyeballed figures. But \textbf{the specification moves
  the answer by 40\%} --- §6.3's misspecification problem in a real
  dataset. Which is why the next step is not optional.
\end{itemize}
\end{frame}

\begin{frame}[fragile]{Step 5: data-driven bandwidths}
\protect\phantomsection\label{step-5-data-driven-bandwidths}
\begin{center}
\small
\begin{tabular}{lccc}
\toprule
 & $p = 0$ & $p = 1$ & $p = 2$ \\
\midrule
Effect of DUI & $-0.0189$ & $-0.0195$ & $-0.0209$ \\
              & $(0.0051)$ & $(0.0061)$ & $(0.0071)$ \\
\midrule
Bandwidth (left)  & $0.020$ & $0.033$ & $0.038$ \\
Bandwidth (right) & $0.014$ & $0.054$ & $0.094$ \\
\bottomrule
\end{tabular}
\end{center}

\begin{itemize}
\tightlist
\item
  \texttt{rdrobust} with a \textbf{triangular} kernel, MSE-optimal
  bandwidths allowed to differ by side, and robust bias-corrected
  inference. This goes beyond what Hansen himself did.
\item
  The estimates sit in a tight band, \(-1.9\) to \(-2.1\) points, and
  are \textbf{far more stable across \(p\)} than the fixed-bandwidth
  regressions were --- and not far from them either.
\item
  Bandwidths \textbf{widen} as \(p\) rises. That is the
  \(n^{-1/(2p+3)}\) formula at work, not a judgement call: more
  polynomial terms mechanically buy a wider window.
\end{itemize}
\end{frame}

\begin{frame}[fragile]{The code}
\protect\phantomsection\label{the-code}
\begin{Shaded}
\begin{Highlighting}[]
\FunctionTok{library}\NormalTok{(rdrobust); }\FunctionTok{library}\NormalTok{(rddensity)}
\NormalTok{d }\OtherTok{\textless{}{-}} \FunctionTok{read.csv}\NormalTok{(}\StringTok{"hansen\_dwi.csv"}\NormalTok{)}

\CommentTok{\# 1. manipulation}
\FunctionTok{rddensity}\NormalTok{(}\AttributeTok{X =}\NormalTok{ d}\SpecialCharTok{$}\NormalTok{bac1, }\AttributeTok{c =} \FloatTok{0.08}\NormalTok{)                       }\CommentTok{\# T = {-}0.14, p = 0.89}

\CommentTok{\# 2. covariate balance, Hansen\textquotesingle{}s window and kernel}
\NormalTok{d}\SpecialCharTok{$}\NormalTok{dui }\OtherTok{\textless{}{-}} \FunctionTok{as.integer}\NormalTok{(d}\SpecialCharTok{$}\NormalTok{bac1 }\SpecialCharTok{\textgreater{}=} \FloatTok{0.08}\NormalTok{); d}\SpecialCharTok{$}\NormalTok{bacc }\OtherTok{\textless{}{-}}\NormalTok{ d}\SpecialCharTok{$}\NormalTok{bac1 }\SpecialCharTok{{-}} \FloatTok{0.08}
\NormalTok{w }\OtherTok{\textless{}{-}} \FunctionTok{subset}\NormalTok{(d, bac1 }\SpecialCharTok{\textgreater{}=} \FloatTok{0.03} \SpecialCharTok{\&}\NormalTok{ bac1 }\SpecialCharTok{\textless{}=} \FloatTok{0.13}\NormalTok{)}
\FunctionTok{lm}\NormalTok{(white }\SpecialCharTok{\textasciitilde{}}\NormalTok{ dui }\SpecialCharTok{*}\NormalTok{ bacc, }\AttributeTok{data =}\NormalTok{ w)                      }\CommentTok{\# repeat for male, aged, acc}

\CommentTok{\# 3. the main regressions}
\FunctionTok{lm}\NormalTok{(recidivism }\SpecialCharTok{\textasciitilde{}}\NormalTok{ dui }\SpecialCharTok{*}\NormalTok{ bacc, }\AttributeTok{data =}\NormalTok{ w)                 }\CommentTok{\# {-}0.0236}
\FunctionTok{lm}\NormalTok{(recidivism }\SpecialCharTok{\textasciitilde{}}\NormalTok{ dui }\SpecialCharTok{*}\NormalTok{ (bacc }\SpecialCharTok{+} \FunctionTok{I}\NormalTok{(bacc}\SpecialCharTok{\^{}}\DecValTok{2}\NormalTok{)), }\AttributeTok{data =}\NormalTok{ w)   }\CommentTok{\# {-}0.0143}
\CommentTok{\# 4. data{-}driven bandwidths}
\FunctionTok{rdrobust}\NormalTok{(}\AttributeTok{y =}\NormalTok{ d}\SpecialCharTok{$}\NormalTok{recidivism, }\AttributeTok{x =}\NormalTok{ d}\SpecialCharTok{$}\NormalTok{bac1, }\AttributeTok{c =} \FloatTok{0.08}\NormalTok{, }\AttributeTok{p =} \DecValTok{1}\NormalTok{,}
         \AttributeTok{kernel =} \StringTok{"triangular"}\NormalTok{, }\AttributeTok{bwselect =} \StringTok{"msetwo"}\NormalTok{)  }\CommentTok{\# {-}0.0195}
\end{Highlighting}
\end{Shaded}

\begin{itemize}
\tightlist
\item
  \texttt{bwselect\ =\ "msetwo"} allows the bandwidth to differ by side.
  With a common bandwidth (\texttt{"mserd"}, the default) the estimate
  is \(-0.0168\): a default is still a choice.
\end{itemize}
\end{frame}

\begin{frame}{What the replication shows}
\protect\phantomsection\label{what-the-replication-shows}
\begin{itemize}
\tightlist
\item
  \textbf{The evidence is coherent.} A continuous density, four flat
  covariates, a visible break in the picture, and estimates between
  \(-1.4\) and \(-2.4\) percentage points across every specification we
  tried.
\item
  \textbf{The steps are the general recipe}, not something specific to
  drunk driving: density of the running variable, covariate balance,
  pictures, parametric regressions, then optimal bandwidths.
\item
  \textbf{The interpretation is less clean than the identification.}
  Getting a DUI charge reduces the chance of another. But is that
  deterrence from the escalated penalty, the shock of the arrest itself,
  or being scared into getting help? RD identifies the effect of
  \emph{crossing the threshold}, not of any single mechanism that
  crossing sets off.
\item
  That gap between a well-identified effect and a well-understood
  mechanism is not a flaw in the design. It is a limit on what any
  threshold can tell you, and it should be stated rather than glossed.
\item
  \textbf{The estimand, precisely:} the average effect of being charged
  with a DUI, for drivers whose BAC was almost exactly \(0.08\). Not for
  drivers at \(0.15\), and not for drivers never stopped.
\end{itemize}
\end{frame}

\section{\texorpdfstring{RD in Finance \emph{(added
section)}}{RD in Finance (added section)}}\label{rd-in-finance-added-section}

\begin{frame}{Where the cutoffs are}
\protect\phantomsection\label{where-the-cutoffs-are}
\begin{itemize}
\tightlist
\item
  Finance is full of rules with thresholds, which is why RD has spread
  quickly in the field.
\end{itemize}

\begin{center}
\small
\begin{tabular}{llll}
\toprule
Paper & Running variable & Cutoff & Design \\
\midrule
Keys et al.\ (2010)       & FICO credit score  & $620$   & sharp-ish \\
Malenko and Shen (2016)   & ISS voting-policy score & policy threshold & sharp \\
Almeida et al.\ (2016)    & EPS relative to consensus & $0$ & sharp \\
Cu\~nat et al.\ (2012)    & shareholder vote share & $50\%$ & fuzzy \\
\bottomrule
\end{tabular}
\end{center}

\begin{itemize}
\tightlist
\item
  Index membership rules, covenant thresholds, credit ratings at the
  investment-grade boundary, regulatory size thresholds --- all
  candidates.
\item
  The econometrics is identical to Hansen's. What differs is how
  plausible \textbf{no manipulation} is: in finance the agents are
  sophisticated, well advised, and highly motivated. Conditions 1--3 of
  §6.4 are routinely all satisfied.
\end{itemize}
\end{frame}

\begin{frame}{Three finance applications}
\protect\phantomsection\label{three-finance-applications}
\begin{itemize}
\tightlist
\item
  \textbf{Keys, Mukherjee, Seru and Vig (2010).} Securitisation rules of
  thumb made loans just above a FICO of \(620\) substantially easier to
  securitise. Loans just above the cutoff default more --- evidence that
  securitisation weakened screening incentives. The running variable is
  a credit score the lender observes and the borrower cannot easily
  move.
\item
  \textbf{Malenko and Shen (2016).} ISS issues a negative say-on-pay
  recommendation when a firm crosses a threshold in its scoring policy.
  Firms just past it see a large drop in voting support, identifying the
  causal influence of proxy advisers rather than merely their
  information.
\item
  \textbf{Almeida, Fos and Kronlund (2016).} Firms that would just miss
  the consensus EPS forecast repurchase shares to push EPS above it.
  Running variable: EPS before repurchases, relative to consensus.
  Cutoff: zero.
\item
  Notice the last one. The whole point of the paper is that firms
  \textbf{manipulate their position} relative to the threshold. That is
  a design where the running variable is endogenous by construction ---
  the assumption violation, studied as the finding.
\end{itemize}
\end{frame}

\begin{frame}{Close votes: the most-used and most-questioned design}
\protect\phantomsection\label{close-votes-the-most-used-and-most-questioned-design}
\begin{itemize}
\tightlist
\item
  A shareholder proposal passing at \(50\%\) looks like the ideal RD.
  Vote shares are near-continuous, the threshold is a legal rule, and
  outcomes are observable. Governance research uses this design heavily.
\item
  \textbf{But are close votes as good as random?} In political science,
  Caughey and Sekhon (2011) found that US House candidates who
  \emph{barely won} differed systematically from those who barely lost
  --- in incumbency, spending, and experience. Density and covariates
  both failed near \(50\%\).
\item
  Eggers, Fowler, Hainmueller, Hall and Snyder (2015) then examined some
  \(40{,}000\) close elections across many settings and found the
  sorting was largely \textbf{specific to the post-war US House}, not a
  general feature of close elections.
\item
  For \textbf{shareholder} votes the concern is sharper than for
  elections. Management sees preliminary tallies, can lobby large
  holders, and can adjourn a meeting. Nudging \(49\%\) to \(51\%\) is
  precisely the sorting McCrary tests for.
\item
  So: run the density test, check covariate balance, and read the
  institutional detail. And note that vote-share designs are usually
  \textbf{fuzzy} --- passing is not implementing.
\end{itemize}
\end{frame}

\begin{frame}{How much data do you actually have?}
\protect\phantomsection\label{how-much-data-do-you-actually-have}
\begin{itemize}
\tightlist
\item
  RD estimates come from a \textbf{neighbourhood}, not a sample. What
  matters is how many observations sit inside the bandwidth, not how
  many rows the dataset has.
\item
  Hansen's data have \(214{,}558\) stops, and the MSE-optimal bandwidths
  still discard most of the range --- with the triangular kernel
  down-weighting even what survives.
\item
  Now consider a typical governance study: a few hundred close votes in
  total, perhaps forty within a few percentage points of \(50\%\). That
  is a very small experiment.
\item
  This is the same idea as the effective sample size from Lecture 3.
  Weighting concentrates influence; RD concentrates it at a point.
\item
  \textbf{Practical consequence:} in finance, RD is often badly
  underpowered, and an underpowered RD produces large, noisy,
  occasionally significant estimates. If you cannot see it in the
  picture, it probably is not there.
\end{itemize}
\end{frame}

\section{6.7 Practical Advice}\label{practical-advice}

\begin{frame}{Visualising the outcome}
\protect\phantomsection\label{visualising-the-outcome}
\begin{itemize}
\tightlist
\item
  RD is an \textbf{unusually visualisation-intensive} method. Get ready
  to make a lot of pictures.
\item
  The rule of thumb: \emph{if you cannot see the effect in the data, you
  are probably underpowered, or it is not there.} A regression that
  finds what the picture does not show deserves suspicion, not a
  footnote.
\item
  What is universal in practice: outcomes plotted as \textbf{means
  within bins} of the running variable. Whether to overlay the fitted
  lines is a matter of taste --- some readers want the raw data alone,
  others want the fit that the regression implies. Both are common.
\item
  Either make the bins yourself, or use automated software.
  \textbf{Binscatter} (Cattaneo, Crump, Farrell and Feng 2024) is the
  popular alternative, and chooses bin widths on a principled basis
  rather than by eye.
\item
  Show \textbf{more than one fit} --- linear and quadratic at least. If
  the discontinuity is visible under both, that is evidence; if it
  appears under only one, that is a specification, not a finding.
\end{itemize}
\end{frame}

\begin{frame}{Show that the answer does not depend on the bandwidth}
\protect\phantomsection\label{show-that-the-answer-does-not-depend-on-the-bandwidth}
\pandocbounded{\includegraphics[keepaspectratio]{04_regression_discontinuity_files/figure-beamer/unnamed-chunk-12-1.pdf}}

\vspace{-0.3em}
\footnotesize

Hansen's data, local linear with a triangular kernel, robust
bias-corrected intervals. The point estimate moves only between
\(-0.016\) and \(-0.021\) across a six-fold range of bandwidths.
\textbf{This is what a robust result looks like.} A plot that swings
across zero, or that is significant only in a narrow window, is telling
you something --- and it is not good news. \hfill \emph{(added)}
\end{frame}

\begin{frame}{Checking the probability of treatment}
\protect\phantomsection\label{checking-the-probability-of-treatment}
\begin{itemize}
\tightlist
\item
  When assignment is \textbf{not} sharp, plot \(\Pr(D = 1)\) against the
  running variable. This is the \textbf{first stage}, and without a
  clear jump --- without ``bite'' --- there is nothing to estimate.
\item
  \textbf{Hoekstra (2009)} is the model. Before estimating the returns
  to attending a state flagship university, he shows that admission at
  the test-score cutoff produces a sharp jump in \emph{attendance}. The
  effect on earnings comes second.
\item
  Two things this picture does for you:

  \begin{enumerate}
  \tightlist
  \item
    It shows the rule is real and enforced, in the data rather than in
    the statute.
  \item
    Its \textbf{size} determines your precision. A first stage of
    \(0.6\) means the reduced form must be scaled up by \(1/0.6\), and
    so must its standard error.
  \end{enumerate}
\item
  Report the first-stage plot even in a design you believe is sharp. If
  the jump is not to one, you have a fuzzy design whether you wanted one
  or not.
\end{itemize}
\end{frame}

\begin{frame}{Checking density and covariate balance}
\protect\phantomsection\label{checking-density-and-covariate-balance}
\begin{itemize}
\tightlist
\item
  \textbf{Density.} Run the McCrary test --- but plot the distribution
  first. A test statistic summarises one feature; the picture shows
  heaping, gaps, mass points, and coarseness that the statistic will
  miss. Almond et al.~is the cautionary case.
\item
  \textbf{Covariate balance.} Plot each covariate against the running
  variable, with the same bins and the same fits as the main result.
  Lee, Moretti and Butler (2004) is the template.
\item
  A jump in a covariate is a real problem: it says the two sides differ
  in something other than treatment, and if an \emph{observed}
  characteristic jumps, unobserved ones probably do too.
\item
  Use the \textbf{same specification} for the balance tests as for the
  main estimate. A balance test at a different bandwidth, or a different
  polynomial order, is not testing the design you actually ran.
\end{itemize}
\end{frame}

\begin{frame}{A checklist for an RD paper}
\protect\phantomsection\label{a-checklist-for-an-rd-paper}
\small

\begin{enumerate}
\tightlist
\item
  \textbf{Name the rule.} Cite the statute, policy document or contract.
  Say who assigns the score and who can see it.
\item
  \textbf{Plot the density} of the running variable and run a formal
  test. Look for heaping, not just a jump.
\item
  \textbf{Plot the outcome} in bins, with fits on each side. If the
  effect is not visible, be sceptical of the regression that finds it.
\item
  \textbf{Show covariate balance}, in a table and a figure, using the
  same specification as the main result.
\item
  \textbf{Report the main estimate} with a triangular kernel, an
  MSE-optimal bandwidth, and \textbf{robust bias-corrected} intervals.
\item
  \textbf{Vary the bandwidth} and plot the estimates. Vary the
  polynomial order too, staying at or below quadratic.
\item
  \textbf{Run placebo cutoffs} where nothing should happen.
\item
  \textbf{State the estimand} as a local effect at the threshold, and
  resist calling it ``the'' effect.
\item
  If the running variable is coarse or heaped, show a
  \textbf{donut-hole} estimate too. \hfill \emph{(added)}
\end{enumerate}
\end{frame}

\section{6.8 Concluding Remarks}\label{concluding-remarks}

\begin{frame}{What the design buys}
\protect\phantomsection\label{what-the-design-buys}
\begin{itemize}
\tightlist
\item
  \textbf{Treatment assignment is known}, not assumed. You can read the
  rule in a statute or a contract, which is what makes RD the most
  credible design among observational methods.
\item
  \textbf{Unconfoundedness is free.} Conditional on the running
  variable, treatment is deterministic --- so the assumption that
  carried all the weight in Lecture 3 is satisfied by construction.
\item
  \textbf{The identifying assumption is about smoothness}, and
  smoothness has testable implications: a continuous density, balanced
  covariates, no effects at placebo cutoffs.
\item
  \textbf{The evidence is visual.} A reader can see the design work, or
  fail, before any regression appears.
\item
  \textbf{And it seems harder to game.} Brodeur, Cook and Heyes (2020)
  find RD studies show less sign of \(p\)-hacking than other
  observational designs --- plausibly because so much raw data is on
  display, the regressions are simple, and there are few researcher
  choices with real power to squeeze water from a rock.
\end{itemize}
\end{frame}

\begin{frame}{What it costs}
\protect\phantomsection\label{what-it-costs}
\begin{itemize}
\tightlist
\item
  \textbf{The estimand is local}, at one point of one variable. It is
  not the \(ATE\) and generally cannot be extrapolated without further
  assumptions.
\item
  \textbf{Overlap is zero}, so everything rests on extrapolating to a
  boundary --- and boundary estimation is exactly where misspecification
  does the most damage. Our cubic example returned \(-170{,}165\)
  against a truth of zero.
\item
  \textbf{The design is data-hungry}, because it estimates effects at
  the boundary of a cutoff. In finance that often means a few dozen
  observations doing the work.
\item
  \textbf{Diagnostics are necessary, not sufficient.} Almond et
  al.~passed the density test, and the donut hole still moved their
  answer by half.
\item
  Acquiring the right data has nothing to do with the econometrics ---
  but if you do not know the econometrics when the opportunity appears,
  you will not recognise it. \textbf{Luck favours the prepared.}
\item
  Next: the fuzzy design, where the cutoff shifts the \emph{probability}
  of treatment --- which makes it an instrument, and takes us to
  instrumental variables.
\end{itemize}
\end{frame}

\begin{frame}{References}
\protect\phantomsection\label{references}
\tiny

\begin{itemize}
\tightlist
\item
  Almeida, H., V. Fos and M. Kronlund (2016). The Real Effects of Share
  Repurchases. \emph{Journal of Financial Economics} 119, 168--185.
\item
  Almond, D., J. Doyle, A. Kowalski and H. Williams (2010). Estimating
  Marginal Returns to Medical Care. \emph{QJE} 125, 591--634.
\item
  Angrist, J. and V. Lavy (1999). Using Maimonides' Rule to Estimate the
  Effect of Class Size. \emph{QJE} 114, 533--575.
\item
  Angrist, J. and M. Rokkanen (2015). Wanna Get Away? RD Estimation Away
  from the Cutoff. \emph{JASA} 110, 1331--1344.
\item
  Barreca, A., M. Guldi, J. Lindo and G. Waddell (2011). Saving Babies?
  \emph{QJE} 126, 2117--2123.
\item
  Becker, G. (1968). Crime and Punishment: An Economic Approach.
  \emph{JPE} 76, 169--217.
\item
  Black, S. (1999). Do Better Schools Matter? \emph{QJE} 114, 577--599.
\item
  Brodeur, A., N. Cook and A. Heyes (2020). Methods Matter. \emph{AER}
  110, 3634--3660.
\item
  Calonico, S., M. Cattaneo and R. Titiunik (2014). Robust Nonparametric
  Confidence Intervals for RD Designs. \emph{Econometrica} 82,
  2295--2326.
\item
  Card, D., C. Dobkin and N. Maestas (2008). The Impact of Nearly
  Universal Insurance Coverage on Health Care Utilization. \emph{AER}
  98, 2242--2258.
\item
  Card, D., C. Dobkin and N. Maestas (2009). Does Medicare Save Lives?
  \emph{QJE} 124, 597--636.
\item
  Carpenter, C. and C. Dobkin (2009). The Effect of Alcohol Consumption
  on Mortality. \emph{AEJ: Applied} 1, 164--182.
\item
  Cattaneo, M., R. Crump, M. Farrell and Y. Feng (2024). On Binscatter.
  \emph{AER} 114, 1488--1514.
\item
  Cattaneo, M., B. Frandsen and R. Titiunik (2015). Randomization
  Inference in the RD Design. \emph{Journal of Causal Inference} 3,
  1--24.
\item
  Cattaneo, M., N. Idrobo and R. Titiunik (2020). \emph{A Practical
  Introduction to Regression Discontinuity Designs}. Cambridge
  University Press.
\item
  Cattaneo, M., M. Jansson and X. Ma (2020). Simple Local Polynomial
  Density Estimators. \emph{JASA} 115, 1449--1455.
\item
  Caughey, D. and J. Sekhon (2011). Elections and the Regression
  Discontinuity Design. \emph{Political Analysis} 19, 385--408.
\item
  Cheng, M.-Y., J. Fan and J. Marron (1997). On Automatic Boundary
  Corrections. \emph{Annals of Statistics} 25, 1691--1708.
\item
  Cunningham, S. \emph{Causal Inference: The Remix}, Ch. 6 --- this
  lecture follows it closely.
\item
  Cu\textasciitilde nat, V., M. Gin\textquotesingle e and M. Guadalupe
  (2012). The Vote Is Cast. \emph{Journal of Finance} 67, 1943--1977.
\item
  Drago, F., R. Galbiati and P. Vertova (2009). The Deterrent Effects of
  Prison. \emph{JPE} 117, 257--280.
\end{itemize}
\end{frame}

\begin{frame}{References (continued)}
\protect\phantomsection\label{references-continued}
\tiny

\begin{itemize}
\tightlist
\item
  Eggers, A., A. Fowler, J. Hainmueller, A. Hall and J. Snyder (2015).
  On the Validity of the RD Design for Electoral Effects. \emph{AJPS}
  59, 259--274.
\item
  Gelman, A. and G. Imbens (2019). Why High-Order Polynomials Should Not
  Be Used in RD Designs. \emph{JBES} 37, 447--456.
\item
  Goldberger, A. (1972). Selection Bias in Evaluating Treatment Effects.
  Discussion paper, Wisconsin.
\item
  Hahn, J., P. Todd and W. van der Klaauw (2001). Identification and
  Estimation with a RD Design. \emph{Econometrica} 69, 201--209.
\item
  Hansen, B. (2015). Punishment and Deterrence: Evidence from Drunk
  Driving. \emph{AER} 105, 1581--1617.
\item
  Hoekstra, M. (2009). The Effect of Attending the Flagship State
  University on Earnings. \emph{ReStat} 91, 717--724.
\item
  Imbens, G. and K. Kalyanaraman (2012). Optimal Bandwidth Choice for
  the RD Estimator. \emph{ReStud} 79, 933--959.
\item
  Imbens, G. and T. Lemieux (2008). Regression Discontinuity Designs: A
  Guide to Practice. \emph{Journal of Econometrics} 142, 615--635.
\item
  Keys, B., T. Mukherjee, A. Seru and V. Vig (2010). Did Securitization
  Lead to Lax Screening? \emph{QJE} 125, 307--362.
\item
  Kolesár, M. and C. Rothe (2018). Inference in RD Designs with a
  Discrete Running Variable. \emph{AER} 108, 2277--2304.
\item
  Lee, D. and J. McCrary (2017). The Deterrence Effect of Prison.
  \emph{Advances in Econometrics} 38, 73--146.
\item
  Lee, D. and T. Lemieux (2010). Regression Discontinuity Designs in
  Economics. \emph{JEL} 48, 281--355.
\item
  Lee, D., E. Moretti and M. Butler (2004). Do Voters Affect or Elect
  Policies? \emph{QJE} 119, 807--859.
\item
  Malenko, N. and Y. Shen (2016). The Role of Proxy Advisory Firms.
  \emph{RFS} 29, 3394--3427.
\item
  McCrary, J. (2008). Manipulation of the Running Variable in the RD
  Design. \emph{Journal of Econometrics} 142, 698--714.
\item
  Thistlethwaite, D. and D. Campbell (1960). Regression-Discontinuity
  Analysis. \emph{Journal of Educational Psychology} 51, 309--317.
\end{itemize}
\end{frame}




\end{document}
